\section{Data and Setting}
\label{sec:data}

\paragraph{The fine-tuning network.}
We use a snapshot of the Hugging Face Hub with metadata and model cards for \nModels{} models by \nAuthorsAll{}
accounts.\todo{Cite the snapshot source and its date.} A model card may declare a base model; the declaration and
its type define a directed edge from the parent (base) to the child. Of the \nEdges{} edges, \nAdapter{} are
adapters and \nFinetune{} full fine-tunes; the remainder are quantized copies and merges. We treat fine-tunes and
adapters (and the few untyped declarations) as \emph{base-model choices}: a developer picked one existing model to
build on. Quantized copies and merges are not choices of this kind; they enter only as attributes of the parent
(e.g., how many quantized versions a candidate has). Creation timestamps are reliable from \tStart{} onward; older
models are kept as part of the initial network, and every child created between \tStart{} and \tEnd{} contributes
one choice per declared parent, \nChoices{} choices in total.

\paragraph{From uploads to adoption decisions.}
Many uploads are repeated variants of one decision: the same author fine-tunes the same base model several times
within hours (checkpoints, hyper-parameter sweeps). We merge consecutive uploads by one author on one parent that
are less than 24 hours apart into a \emph{decision} (\nDecisions{} decisions; 6- and 72-hour windows in
\cref{app:repeat}). Our main sample is the first decision of each author--parent pair, a \emph{first adoption}:
\nFirst{} first adoptions by \nFirstAuthors{} authors of \nFirstParents{} distinct parents. First adoptions are the
decisions in which a developer selects a base model rather than continuing to use one; repeated use is a different
behavior, which we analyze separately in \cref{app:repeat}. The \emph{market} for a first adoption is every model of
the same modality (text, image generation, speech, vision, multimodal) created before the child, minus the parents
the author has already used. Markets are large: the median first adoption faces \marketMedian{} alternatives, and
\shareText\% of first adoptions are text models.

\paragraph{A concentrated network.}
Adoption is extremely concentrated. The single most adopted parent receives \topOneReal\% of all first adoptions,
the ten most adopted receive \topTenReal\%, and the effective number of parents, $\exp(H)$ for the entropy $H$ of
the adoption shares, is \effReal{} out of \nFirstParents{} parents ever adopted. Concentration is what popularity
would produce, but also what a few models of exceptional general quality, or a few models that fit the most common
needs, would produce. Separating these explanations is the purpose of the choice model.

\paragraph{What a developer can see: four groups of signals.}
For every (choice, candidate) pair we measure signals available to the developer at the time of the choice, from the
network state and metadata strictly before the child's creation. We organize them into four groups
(\cref{app:data} lists every feature):
\begin{itemize}
\item \textbf{Popularity}: signals that are the same for every developer and do not depend on what the developer
needs: degree $d$ and $\1[d>0]$, adoptions in the last 30 days, age, a major-organization publisher, a legacy
model (created before March 2022), the library ecosystem, the number of quantized versions, re-distributor accounts that republish
other organizations' weights, co-use by other developers who trained on the same datasets or used the same niche
parents, and whether size, card text or benchmark scores are available at all (documentation completeness is
visibility, not demand).
\item \textbf{General quality}: model-index benchmark scores reported in the card, standardized within benchmark
(a z-score). Hub fixed effects, which absorb any unobserved quality, are used only to bound this group
(\cref{sec:bounds}).
\item \textbf{Demand match}: how well the candidate fits this developer's need. Need is read only from sources the
child does not copy from its parent: the author's earlier parents (task, languages, size; the candidate itself left
out), the languages and tasks common among other models trained on the child's declared datasets, and the TF-IDF
similarity between the child's and the candidate's model cards.
\item \textbf{Inertia}: the candidate is the author's own model, or belongs to a model family the author used before.
\end{itemize}

\paragraph{Guarding against copied metadata.}
A child's card is partly inherited from its parent, so naive similarity between child and candidate would favor the
true parent mechanically. We test each field for inheritance by comparing the child with its parent and with a
random parent from the same month and modality (Cohen's $\kappa$, \cref{app:data}). Parameter count ($\kappa=0.84$)
and license ($\kappa=0.74$) are mostly copied and are not used as demand; datasets ($\kappa=0.12$) are mostly the
child's own. For dataset-based demand we additionally drop datasets the child shares with its chosen parent, compute
each dataset's language and task profile only from models created before the choice, and remove the candidate's own
children from the profile used for that candidate. For card text we drop template paragraphs (paragraphs in
$\ge$20 cards, including near-duplicates found by MinHash), paragraphs copied from an earlier card, and organization
and family names, so that a shared model name cannot produce similarity.

\paragraph{Training and holdout.}
We hold out the last 10\% of choices in time (children created on or after \tHold) for out-of-sample evaluation;
\nTrain{} first adoptions are used for estimation. The fixed effects of \cref{sec:bounds} are assigned to the 500
most adopted parents as of the holdout cut, so the set does not depend on future adoptions.
